\documentclass[11pt,a4paper]{report}

% Nombre y versión de este manual: los usan la portada, el encabezado y el PDF.
\newcommand{\nombremanual}{Manual de Usuario}
\newcommand{\versionmanual}{noviembre de 2026}

% Paquetes, paleta, títulos, cajas y portada, compartidos con el Manual de Administración.
% ======================================================================
% Estilo común de los manuales del Sistema de Biblioteca UMAG.
% Lo cargan manual-usuario.tex y manual-admin.tex con % ======================================================================
% Estilo común de los manuales del Sistema de Biblioteca UMAG.
% Lo cargan manual-usuario.tex y manual-admin.tex con % ======================================================================
% Estilo común de los manuales del Sistema de Biblioteca UMAG.
% Lo cargan manual-usuario.tex y manual-admin.tex con \input{estilo-manual},
% después de definir \nombremanual y \versionmanual.
% ======================================================================

% ---------- Codificación e idioma ----------
\usepackage[utf8]{inputenc}
\usepackage[T1]{fontenc}
\usepackage[spanish,es-noindentfirst]{babel}

% ---------- Diseño de página ----------
\usepackage[a4paper,margin=2.5cm,headheight=15pt]{geometry}
\usepackage{xcolor}
\usepackage{graphicx}
\usepackage{fancyhdr}
\usepackage{titlesec}
\usepackage{enumitem}
\usepackage[most]{tcolorbox}
\usepackage{booktabs}
\usepackage{tabularx}
\usepackage{caption}
\usepackage{needspace}
\usepackage{etoolbox}
\usepackage{amssymb}
\usepackage{hyperref}

% ---------- Paleta institucional (la misma de la app: docs/estilo-visual.md) ----------
% Morados del degradado de la barra superior y dorado de la pestaña activa.
\definecolor{umagpurple}{HTML}{3B28A3}
\definecolor{umagpurpledark}{HTML}{2D1B69}
\definecolor{umaggold}{HTML}{B08D57}
\definecolor{umaggray}{HTML}{4B5563}
\definecolor{cajafondo}{HTML}{F8FAFC}

\hypersetup{
  colorlinks=true,
  linkcolor=umagpurple,
  urlcolor=umagpurple,
  citecolor=umagpurple,
  pdftitle={\nombremanual{} - Sistema de Biblioteca UMAG},
  pdfauthor={Biblioteca UMAG},
}

% ---------- Encabezados / pies de página ----------
\pagestyle{fancy}
\fancyhf{}
\renewcommand{\headrulewidth}{0.4pt}
\renewcommand{\footrulewidth}{0pt}
% Encabezado de capítulo sin el prefijo "Capítulo N." (ya se ve grande en la
% portadilla del capítulo) para que quepa junto al nombre del manual sin superponerse.
\renewcommand{\chaptermark}[1]{\markboth{#1}{}}
\fancyhead[L]{\small\textcolor{umaggray}{Biblioteca UMAG · \nombremanual}}
\fancyhead[R]{\small\textcolor{umaggray}{\leftmark}}
\fancyfoot[C]{\small\textcolor{umaggray}{\thepage}}

\fancypagestyle{plain}{%
  \fancyhf{}
  \fancyfoot[C]{\small\textcolor{umaggray}{\thepage}}
  \renewcommand{\headrulewidth}{0pt}
}

% ---------- Estilo de títulos ----------
\titleformat{\chapter}[display]
  {\normalfont\bfseries\color{umagpurple}}
  {\Large\filright\textcolor{umaggold}{CAPÍTULO \thechapter}}
  {6pt}
  {\Huge\filright}
\titlespacing*{\chapter}{0pt}{0pt}{28pt}

\titleformat{\section}
  {\normalfont\Large\bfseries\color{umagpurpledark}}{\thesection}{0.6em}{}
\titleformat{\subsection}
  {\normalfont\large\bfseries\color{umagpurple}}{\thesubsection}{0.6em}{}
% Un título de sección no queda solo al pie de la página: si no caben él,
% su etiqueta de permiso y unas líneas de texto, pasa a la página siguiente.
\pretocmd{\section}{\needspace{7\baselineskip}}{}{}

% ---------- Cajas reutilizables ----------
\newtcolorbox{nota}{
  colback=umaggold!8, colframe=umaggold, boxrule=0.6pt, arc=2pt,
  left=8pt, right=8pt, top=6pt, bottom=6pt,
  fonttitle=\bfseries, title={Tip},
}
\newtcolorbox{importante}{
  colback=umagpurple!6, colframe=umagpurple, boxrule=0.6pt, arc=2pt,
  left=8pt, right=8pt, top=6pt, bottom=6pt,
  fonttitle=\bfseries, title={Importante},
}
% Etiqueta compacta (no una caja que envuelve contenido) para marcar una
% sección que solo ven los roles con cierto permiso, sin arriesgar que una
% tcolorbox con una captura de pantalla adentro desborde el alto de la página.
% \requierepermiso{nombre del permiso, tal como aparece en Personal y Roles}
\newtcbox{\etiquetapermiso}{on line, colback=umagpurple, colframe=umagpurple,
  coltext=white, boxrule=0pt, arc=0pt, left=4pt, right=4pt, top=2pt, bottom=2pt,
  fontupper=\bfseries\small}
\newcommand{\requierepermiso}[1]{%
  \needspace{3\baselineskip}%
  \par\noindent\etiquetapermiso{REQUIERE PERMISO: #1}\par\vspace{6pt}%
}

% ---------- Comando para insertar capturas de pantalla ----------
% \captura{archivo}{texto de pie de figura}
\newcommand{\captura}[2]{%
  \begin{center}
  \fbox{\includegraphics[width=0.92\textwidth,height=0.78\textheight,keepaspectratio]{imagenes/#1}}\\[5pt]
  {\small\itshape\color{umaggray} #2}
  \end{center}
  \vspace{6pt}
}

\setlist[itemize]{leftmargin=1.4em}
\setlist[enumerate]{leftmargin=1.6em}

% ---------- Portada ----------
% \portadamanual{subtítulo}{contenido del recuadro de acceso}
\newcommand{\portadamanual}[2]{%
\begin{titlepage}
\centering
\vspace*{1.5cm}

{\includegraphics[width=3cm]{imagenes/logo-umag.png}}

\vspace{0.8cm}

{\color{umaggray}\large UNIVERSIDAD DE MAGALLANES}

\vspace{2.2cm}

{\color{umagpurple}\Huge\bfseries Sistema de Biblioteca}\\[6pt]
{\color{umagpurpledark}\Huge\bfseries UMAG}

\vspace{1cm}
{\color{umaggold}\rule{8cm}{1.2pt}}
\vspace{1cm}

{\color{umaggray}\LARGE \nombremanual}

\vspace{0.4cm}
{\color{umaggray}\large #1}

\vfill

\begin{tcolorbox}[colback=cajafondo,colframe=umagpurple!40,width=0.8\textwidth,arc=3pt,boxrule=0.6pt]
\centering
#2
\end{tcolorbox}

\vspace{1cm}
{\color{umaggray} Versión del documento: \versionmanual}

\vspace{1cm}
\end{titlepage}
}
,
% después de definir \nombremanual y \versionmanual.
% ======================================================================

% ---------- Codificación e idioma ----------
\usepackage[utf8]{inputenc}
\usepackage[T1]{fontenc}
\usepackage[spanish,es-noindentfirst]{babel}

% ---------- Diseño de página ----------
\usepackage[a4paper,margin=2.5cm,headheight=15pt]{geometry}
\usepackage{xcolor}
\usepackage{graphicx}
\usepackage{fancyhdr}
\usepackage{titlesec}
\usepackage{enumitem}
\usepackage[most]{tcolorbox}
\usepackage{booktabs}
\usepackage{tabularx}
\usepackage{caption}
\usepackage{needspace}
\usepackage{etoolbox}
\usepackage{amssymb}
\usepackage{hyperref}

% ---------- Paleta institucional (la misma de la app: docs/estilo-visual.md) ----------
% Morados del degradado de la barra superior y dorado de la pestaña activa.
\definecolor{umagpurple}{HTML}{3B28A3}
\definecolor{umagpurpledark}{HTML}{2D1B69}
\definecolor{umaggold}{HTML}{B08D57}
\definecolor{umaggray}{HTML}{4B5563}
\definecolor{cajafondo}{HTML}{F8FAFC}

\hypersetup{
  colorlinks=true,
  linkcolor=umagpurple,
  urlcolor=umagpurple,
  citecolor=umagpurple,
  pdftitle={\nombremanual{} - Sistema de Biblioteca UMAG},
  pdfauthor={Biblioteca UMAG},
}

% ---------- Encabezados / pies de página ----------
\pagestyle{fancy}
\fancyhf{}
\renewcommand{\headrulewidth}{0.4pt}
\renewcommand{\footrulewidth}{0pt}
% Encabezado de capítulo sin el prefijo "Capítulo N." (ya se ve grande en la
% portadilla del capítulo) para que quepa junto al nombre del manual sin superponerse.
\renewcommand{\chaptermark}[1]{\markboth{#1}{}}
\fancyhead[L]{\small\textcolor{umaggray}{Biblioteca UMAG · \nombremanual}}
\fancyhead[R]{\small\textcolor{umaggray}{\leftmark}}
\fancyfoot[C]{\small\textcolor{umaggray}{\thepage}}

\fancypagestyle{plain}{%
  \fancyhf{}
  \fancyfoot[C]{\small\textcolor{umaggray}{\thepage}}
  \renewcommand{\headrulewidth}{0pt}
}

% ---------- Estilo de títulos ----------
\titleformat{\chapter}[display]
  {\normalfont\bfseries\color{umagpurple}}
  {\Large\filright\textcolor{umaggold}{CAPÍTULO \thechapter}}
  {6pt}
  {\Huge\filright}
\titlespacing*{\chapter}{0pt}{0pt}{28pt}

\titleformat{\section}
  {\normalfont\Large\bfseries\color{umagpurpledark}}{\thesection}{0.6em}{}
\titleformat{\subsection}
  {\normalfont\large\bfseries\color{umagpurple}}{\thesubsection}{0.6em}{}
% Un título de sección no queda solo al pie de la página: si no caben él,
% su etiqueta de permiso y unas líneas de texto, pasa a la página siguiente.
\pretocmd{\section}{\needspace{7\baselineskip}}{}{}

% ---------- Cajas reutilizables ----------
\newtcolorbox{nota}{
  colback=umaggold!8, colframe=umaggold, boxrule=0.6pt, arc=2pt,
  left=8pt, right=8pt, top=6pt, bottom=6pt,
  fonttitle=\bfseries, title={Tip},
}
\newtcolorbox{importante}{
  colback=umagpurple!6, colframe=umagpurple, boxrule=0.6pt, arc=2pt,
  left=8pt, right=8pt, top=6pt, bottom=6pt,
  fonttitle=\bfseries, title={Importante},
}
% Etiqueta compacta (no una caja que envuelve contenido) para marcar una
% sección que solo ven los roles con cierto permiso, sin arriesgar que una
% tcolorbox con una captura de pantalla adentro desborde el alto de la página.
% \requierepermiso{nombre del permiso, tal como aparece en Personal y Roles}
\newtcbox{\etiquetapermiso}{on line, colback=umagpurple, colframe=umagpurple,
  coltext=white, boxrule=0pt, arc=0pt, left=4pt, right=4pt, top=2pt, bottom=2pt,
  fontupper=\bfseries\small}
\newcommand{\requierepermiso}[1]{%
  \needspace{3\baselineskip}%
  \par\noindent\etiquetapermiso{REQUIERE PERMISO: #1}\par\vspace{6pt}%
}

% ---------- Comando para insertar capturas de pantalla ----------
% \captura{archivo}{texto de pie de figura}
\newcommand{\captura}[2]{%
  \begin{center}
  \fbox{\includegraphics[width=0.92\textwidth,height=0.78\textheight,keepaspectratio]{imagenes/#1}}\\[5pt]
  {\small\itshape\color{umaggray} #2}
  \end{center}
  \vspace{6pt}
}

\setlist[itemize]{leftmargin=1.4em}
\setlist[enumerate]{leftmargin=1.6em}

% ---------- Portada ----------
% \portadamanual{subtítulo}{contenido del recuadro de acceso}
\newcommand{\portadamanual}[2]{%
\begin{titlepage}
\centering
\vspace*{1.5cm}

{\includegraphics[width=3cm]{imagenes/logo-umag.png}}

\vspace{0.8cm}

{\color{umaggray}\large UNIVERSIDAD DE MAGALLANES}

\vspace{2.2cm}

{\color{umagpurple}\Huge\bfseries Sistema de Biblioteca}\\[6pt]
{\color{umagpurpledark}\Huge\bfseries UMAG}

\vspace{1cm}
{\color{umaggold}\rule{8cm}{1.2pt}}
\vspace{1cm}

{\color{umaggray}\LARGE \nombremanual}

\vspace{0.4cm}
{\color{umaggray}\large #1}

\vfill

\begin{tcolorbox}[colback=cajafondo,colframe=umagpurple!40,width=0.8\textwidth,arc=3pt,boxrule=0.6pt]
\centering
#2
\end{tcolorbox}

\vspace{1cm}
{\color{umaggray} Versión del documento: \versionmanual}

\vspace{1cm}
\end{titlepage}
}
,
% después de definir \nombremanual y \versionmanual.
% ======================================================================

% ---------- Codificación e idioma ----------
\usepackage[utf8]{inputenc}
\usepackage[T1]{fontenc}
\usepackage[spanish,es-noindentfirst]{babel}

% ---------- Diseño de página ----------
\usepackage[a4paper,margin=2.5cm,headheight=15pt]{geometry}
\usepackage{xcolor}
\usepackage{graphicx}
\usepackage{fancyhdr}
\usepackage{titlesec}
\usepackage{enumitem}
\usepackage[most]{tcolorbox}
\usepackage{booktabs}
\usepackage{tabularx}
\usepackage{caption}
\usepackage{needspace}
\usepackage{etoolbox}
\usepackage{amssymb}
\usepackage{hyperref}

% ---------- Paleta institucional (la misma de la app: docs/estilo-visual.md) ----------
% Morados del degradado de la barra superior y dorado de la pestaña activa.
\definecolor{umagpurple}{HTML}{3B28A3}
\definecolor{umagpurpledark}{HTML}{2D1B69}
\definecolor{umaggold}{HTML}{B08D57}
\definecolor{umaggray}{HTML}{4B5563}
\definecolor{cajafondo}{HTML}{F8FAFC}

\hypersetup{
  colorlinks=true,
  linkcolor=umagpurple,
  urlcolor=umagpurple,
  citecolor=umagpurple,
  pdftitle={\nombremanual{} - Sistema de Biblioteca UMAG},
  pdfauthor={Biblioteca UMAG},
}

% ---------- Encabezados / pies de página ----------
\pagestyle{fancy}
\fancyhf{}
\renewcommand{\headrulewidth}{0.4pt}
\renewcommand{\footrulewidth}{0pt}
% Encabezado de capítulo sin el prefijo "Capítulo N." (ya se ve grande en la
% portadilla del capítulo) para que quepa junto al nombre del manual sin superponerse.
\renewcommand{\chaptermark}[1]{\markboth{#1}{}}
\fancyhead[L]{\small\textcolor{umaggray}{Biblioteca UMAG · \nombremanual}}
\fancyhead[R]{\small\textcolor{umaggray}{\leftmark}}
\fancyfoot[C]{\small\textcolor{umaggray}{\thepage}}

\fancypagestyle{plain}{%
  \fancyhf{}
  \fancyfoot[C]{\small\textcolor{umaggray}{\thepage}}
  \renewcommand{\headrulewidth}{0pt}
}

% ---------- Estilo de títulos ----------
\titleformat{\chapter}[display]
  {\normalfont\bfseries\color{umagpurple}}
  {\Large\filright\textcolor{umaggold}{CAPÍTULO \thechapter}}
  {6pt}
  {\Huge\filright}
\titlespacing*{\chapter}{0pt}{0pt}{28pt}

\titleformat{\section}
  {\normalfont\Large\bfseries\color{umagpurpledark}}{\thesection}{0.6em}{}
\titleformat{\subsection}
  {\normalfont\large\bfseries\color{umagpurple}}{\thesubsection}{0.6em}{}
% Un título de sección no queda solo al pie de la página: si no caben él,
% su etiqueta de permiso y unas líneas de texto, pasa a la página siguiente.
\pretocmd{\section}{\needspace{7\baselineskip}}{}{}

% ---------- Cajas reutilizables ----------
\newtcolorbox{nota}{
  colback=umaggold!8, colframe=umaggold, boxrule=0.6pt, arc=2pt,
  left=8pt, right=8pt, top=6pt, bottom=6pt,
  fonttitle=\bfseries, title={Tip},
}
\newtcolorbox{importante}{
  colback=umagpurple!6, colframe=umagpurple, boxrule=0.6pt, arc=2pt,
  left=8pt, right=8pt, top=6pt, bottom=6pt,
  fonttitle=\bfseries, title={Importante},
}
% Etiqueta compacta (no una caja que envuelve contenido) para marcar una
% sección que solo ven los roles con cierto permiso, sin arriesgar que una
% tcolorbox con una captura de pantalla adentro desborde el alto de la página.
% \requierepermiso{nombre del permiso, tal como aparece en Personal y Roles}
\newtcbox{\etiquetapermiso}{on line, colback=umagpurple, colframe=umagpurple,
  coltext=white, boxrule=0pt, arc=0pt, left=4pt, right=4pt, top=2pt, bottom=2pt,
  fontupper=\bfseries\small}
\newcommand{\requierepermiso}[1]{%
  \needspace{3\baselineskip}%
  \par\noindent\etiquetapermiso{REQUIERE PERMISO: #1}\par\vspace{6pt}%
}

% ---------- Comando para insertar capturas de pantalla ----------
% \captura{archivo}{texto de pie de figura}
\newcommand{\captura}[2]{%
  \begin{center}
  \fbox{\includegraphics[width=0.92\textwidth,height=0.78\textheight,keepaspectratio]{imagenes/#1}}\\[5pt]
  {\small\itshape\color{umaggray} #2}
  \end{center}
  \vspace{6pt}
}

\setlist[itemize]{leftmargin=1.4em}
\setlist[enumerate]{leftmargin=1.6em}

% ---------- Portada ----------
% \portadamanual{subtítulo}{contenido del recuadro de acceso}
\newcommand{\portadamanual}[2]{%
\begin{titlepage}
\centering
\vspace*{1.5cm}

{\includegraphics[width=3cm]{imagenes/logo-umag.png}}

\vspace{0.8cm}

{\color{umaggray}\large UNIVERSIDAD DE MAGALLANES}

\vspace{2.2cm}

{\color{umagpurple}\Huge\bfseries Sistema de Biblioteca}\\[6pt]
{\color{umagpurpledark}\Huge\bfseries UMAG}

\vspace{1cm}
{\color{umaggold}\rule{8cm}{1.2pt}}
\vspace{1cm}

{\color{umaggray}\LARGE \nombremanual}

\vspace{0.4cm}
{\color{umaggray}\large #1}

\vfill

\begin{tcolorbox}[colback=cajafondo,colframe=umagpurple!40,width=0.8\textwidth,arc=3pt,boxrule=0.6pt]
\centering
#2
\end{tcolorbox}

\vspace{1cm}
{\color{umaggray} Versión del documento: \versionmanual}

\vspace{1cm}
\end{titlepage}
}


\begin{document}

\portadamanual
  {Portal de la biblioteca para estudiantes,\\docentes y funcionarios}
  {\textbf{Portal de la biblioteca:} ruta \texttt{/login}, con tu correo institucional\\[6pt]
   {\small (entra desde la dirección web que indique la biblioteca)}}

\tableofcontents
\clearpage

% ======================================================================
\chapter{Introducción}
% ======================================================================

\section{Sobre este manual}

Este manual explica, paso a paso, cómo usar el \textbf{portal de la
biblioteca}: el sitio de autoservicio del Sistema de Biblioteca UMAG para
estudiantes, docentes y funcionarios. Desde el portal puedes, sin pasar por
el mesón:

\begin{itemize}
  \item marcar tu entrada a la biblioteca, escaneando un código QR o con tu
  RUT;
  \item buscar libros y revistas en el catálogo, ver si están disponibles y
  reservarlos;
  \item reservar una sala de estudio para el día de hoy;
  \item revisar tus préstamos, reservas y multas, y descargar tu Constancia
  de No Multa;
  \item ver tu historial de ingresos.
\end{itemize}

Llevarte un libro o un recurso de apoyo, devolverlo y pagar una multa se
hace en el mesón de la biblioteca; el capítulo~\ref{cap:meson} explica las
reglas.

Cada capítulo corresponde a una pantalla del portal, con capturas reales y
una explicación de qué hace cada botón.

\section{Requisitos}

\begin{itemize}
  \item Un computador, tablet o teléfono con un navegador web moderno
  (Chrome, Edge, Firefox).
  \item Conexión a la red donde esté publicado el sistema.
  \item Tu correo institucional y tu contraseña del portal
  (capítulo~\ref{cap:login}).
  \item Para marcar tu entrada con el código QR, un teléfono con cámara.
\end{itemize}

\section{Tu cuenta}

Tu cuenta de la biblioteca está \textbf{vigente hasta el 31 de marzo} del
año siguiente a su última renovación. Con la cuenta vencida no puedes
marcar tu entrada, entrar al portal ni llevarte material, pero sí devolver
lo que tengas y pagar tus multas. Para renovarla, acércate al mesón.

\section{Cómo moverse por el portal}

El portal se adapta al equipo desde el que lo abres. En un computador
(ventanas de 1024 píxeles de ancho o más) muestra un panel de inicio con tu
actividad, las dos formas de marcar entrada lado a lado y el catálogo junto
a \textbf{Mis reservas}. En teléfonos y tablets usa una sola columna, con
botones grandes pensados para el tacto.

En la barra superior, las pestañas \textbf{Inicio}, \textbf{Entrada},
\textbf{Catálogo} y \textbf{Salas} llevan a cada sección. A la derecha
aparecen tu nombre y tu tipo de usuario; al pulsarlos se abre un menú con
el RUT de tu cuenta, \textbf{Mi perfil}, \textbf{Mis ingresos} y
\textbf{Cerrar sesión}. En el teléfono, una barra fija abajo lleva a
Inicio, Catálogo, Entrada (el botón central), Salas y Perfil.

Junto a tu nombre está el interruptor \textbf{día / noche}: al pulsarlo, el
portal pasa a colores oscuros, más descansados para la vista de noche, y al
volver a pulsarlo regresa al modo claro. La elección queda guardada en ese
teléfono o computador.

Las capturas de este manual muestran el formato de computador; el de
teléfono se ve así:

\captura{27b-portal-telefono.png}{El portal en un teléfono: inicio, marcar entrada, catálogo y salas (la línea de tiempo se desliza de lado).}

% ======================================================================
\chapter{Inicio de sesión}
\label{cap:login}
% ======================================================================

Al portal se entra por el mismo inicio de sesión que usa el personal de la
biblioteca, en la ruta \texttt{/login}. El sistema reconoce que tu cuenta
es de estudiante, docente o funcionario y te lleva al portal.

\captura{01-login.png}{Inicio de sesión, el mismo para estudiantes, docentes, funcionarios y personal de biblioteca.}

\begin{enumerate}
  \item En \textbf{Correo institucional}, escribe tu correo de la
  universidad (por ejemplo, \texttt{nombre@umag.cl}). Si la biblioteca
  tiene registrado otro correo tuyo, también sirve.
  \item En \textbf{Contraseña}, escribe tu contraseña del portal. El ojo al
  final de la casilla la muestra u oculta.
  \item Pulsa \textbf{Ingresar} (o \texttt{Enter}).
\end{enumerate}

\begin{importante}
Se entra \textbf{con el correo, no con el RUT}. Si escribes tu RUT, el
sistema avisa: ``Se entra con el correo institucional (por ejemplo,
nombre@umag.cl), no con el RUT.'' El RUT se sigue usando para marcar tu
entrada y en las reservas de salas.
\end{importante}

\captura{01b-login-con-rut.png}{Con el RUT no se entra: el aviso pide el correo institucional.}

Si llegaste al inicio de sesión desde otra página ---por ejemplo, al
escanear el código QR de la entrada sin haber iniciado sesión---, al entrar
vuelves a esa misma página.

\section{Si también tienes cuenta del personal}

Quien trabaja en la biblioteca y además estudia puede tener las dos cuentas
con el mismo correo. Si las dos tienen la misma contraseña, después de
pulsar \textbf{Ingresar} aparece \textbf{Tienes dos cuentas}: elige
\textbf{Portal de la biblioteca} para entrar al portal, o \textbf{Usar
otros datos} para volver al formulario. Si cada cuenta tiene su propia
contraseña, entras a la cuenta de la contraseña que escribiste.

\captura{01c-login-dos-cuentas.png}{Una persona con cuenta del personal y de estudiante elige con cuál entra.}

\section{Si olvidaste tu contraseña}

El portal no tiene recuperación de contraseña en línea: el enlace
\textbf{¿Olvidaste tu contraseña?} recuerda que hay que acercarse al mesón
de la biblioteca (o contactar a un administrador). Lleva tu credencial y el
personal te asigna una nueva.

\begin{nota}
Los botones \textbf{Continuar con Google} y \textbf{Cuenta institucional
(LDAP)} servirán para entrar con tu cuenta de la universidad cuando esta
los habilite. Mientras tanto, se entra con el correo y la contraseña.
\end{nota}

% ======================================================================
\chapter{Página de Inicio}
% ======================================================================

En el computador, el inicio es un panel con lo más importante del día:

\captura{27-portal-home.png}{Inicio del portal en el computador: ingreso de hoy, actividad y accesos rápidos.}

\begin{itemize}
  \item \textbf{¿Llegaste a la biblioteca?:} indica si ya marcaste tu
  ingreso hoy (y a qué hora) y ofrece \textbf{Escanear QR} o \textbf{Usar
  mi RUT}, que llevan directo a Marcar Entrada con la cámara encendida o
  con el cursor en el campo del RUT.
  \item \textbf{Ingresos de hoy:} cuántas personas marcaron su entrada hoy,
  frente al aforo de la biblioteca. Como el sistema no registra salidas, no
  es la cantidad de personas que están adentro en este momento.
  \item \textbf{Buscador:} busca en el catálogo por título, autor o ISBN.
  \item \textbf{Tu actividad:} tus préstamos vigentes con su fecha de
  devolución (con un aviso si vencen en los próximos días o están
  atrasados), tus reservas de libros, tus reservas de sala de hoy y el
  estado de tus multas. Sin multas pendientes, desde ahí descargas la
  \textbf{Constancia de No Multa} en PDF.
  \item \textbf{Accesos rápidos:} reservar sala, explorar el catálogo y
  \textbf{Recursos digitales} (\url{https://umag.elogim.com/}), la base de
  datos digital de la universidad, que es un sitio externo al sistema de
  biblioteca.
\end{itemize}

Arriba a la derecha, una etiqueta indica si la biblioteca está abierta
(«Abierta · cierra a las 20:30») o, si está cerrada, cuándo vuelve a abrir
(por ejemplo, «Cerrada · abre el lunes a las 08:45»), según el horario de
atención: lunes a viernes de 08:45 a 20:30 y sábado de 09:00 a 14:00.

En el teléfono, el inicio saluda con la fecha, el horario de la biblioteca
y los ingresos del día. Debajo están el aviso \textbf{¿Llegaste a la
biblioteca?} con \textbf{Escanear QR} y \textbf{Usar mi RUT}, el buscador
del catálogo y \textbf{Explorar}: Catálogo, Reservar sala, Recursos
digitales y Mis ingresos. La Constancia de No Multa, en el teléfono, está
en \textbf{Perfil}.

% ======================================================================
\chapter{Marcar tu entrada}
% ======================================================================

Al llegar a la biblioteca puedes registrar tu propio ingreso de dos formas,
cada una en su tarjeta. Arriba a la derecha se indica si la biblioteca está
abierta y hasta qué hora.

\captura{28-portal-entrada.png}{Marcar Entrada: escanear el código QR o ingresar con el RUT.}

\begin{itemize}
  \item \textbf{Escanear código QR} (recomendado en el celular):
  \textbf{Abrir cámara} enciende la cámara dentro de la misma tarjeta; al
  leer el código QR de la biblioteca, tu entrada queda registrada. La
  primera vez, el navegador pide permiso para usar la cámara.
  \item \textbf{Ingresar con RUT} (ideal en el computador): escribe tu RUT
  sin puntos ni guion ---el sistema lo formatea--- y pulsa
  \textbf{Registrar entrada} o \textbf{Enter}.
\end{itemize}

Si algo falla ---por ejemplo, un RUT que no es el de tu cuenta--- el aviso
aparece en la misma tarjeta, sin salir de la pantalla.

\captura{28b-portal-entrada-rut.png}{El RUT escrito no es el de la cuenta: el aviso aparece en la tarjeta.}

\begin{nota}
El RUT que escribes tiene que ser exactamente el de tu cuenta: desde tu
sesión no se puede marcar la entrada de otra persona.
\end{nota}

Con el RUT correcto, tu entrada queda registrada de inmediato y el sistema
muestra una bienvenida antes de volver solo al inicio del portal.

\captura{28c-portal-entrada-exito.png}{Confirmación de entrada registrada.}

% ======================================================================
\chapter{Catálogo y reservas de libros}
% ======================================================================

En el catálogo consultas la disponibilidad de libros y revistas. El
buscador acepta título, autor, área o ISBN (con o sin guiones) e indica
cuántos resultados hay; los botones \textbf{Todos} y \textbf{Disponibles
ahora} y el selector de \textbf{Área} filtran la lista. Al pulsar
\textbf{Área} aparecen las áreas con más títulos; para encontrar otra,
escribe unas letras de su nombre y elígela de la lista (la cruz $\times$
quita el filtro).

\captura{29-portal-catalogo.png}{Catálogo del portal, con filtros, disponibilidad y reserva directa.}

Cada título muestra una portada con sus iniciales (del color de su área),
su estado y la acción que corresponde:

\begin{itemize}
  \item \textbf{Disponible}, con cuántos ejemplares hay en estantería:
  botón \textbf{Reservar}.
  \item \textbf{Prestado}, con la fecha aproximada en que vuelve una copia:
  botón \textbf{Avisarme al volver}, que te deja en la lista de espera.
  \textbf{No disponible} indica que no hay copias libres aunque ninguna
  esté prestada (por ejemplo, porque está apartada para otra persona).
  \item \textbf{Lista para retirar} o \textbf{En lista de espera}, si ya lo
  reservaste: en lugar del botón aparece \textbf{En tus reservas}.
\end{itemize}

\textbf{Ver ficha}, junto a cada título, muestra sus datos bibliográficos
(ISBN, clasificación, pie de imprenta, materias, acceso
electrónico\ldots). En el computador se despliega dentro de la lista; en el
teléfono se abre una ventana con el resumen del título ---portada, estado,
ficha y el botón para reservar--- sin alargar la página.

\captura{29b-portal-ficha.png}{Ficha de un título en el portal.}

\section{Reservar un libro}

Reservar siempre pide confirmación. \textbf{Reservar} y \textbf{Avisarme
al volver} abren una ventana con el título y lo que va a pasar:

\begin{itemize}
  \item Con \textbf{Reservar} se aparta una copia a tu nombre, que retiras
  en el mesón antes de la fecha límite: tienes \textbf{tres días}, y si el
  plazo cae en un día en que la biblioteca no atiende, se corre al
  siguiente.
  \item Con \textbf{Avisarme al volver} quedas en la lista de espera.
  Cuando llega tu turno se te aparta una copia y se te avisa por correo, si
  tienes uno registrado.
\end{itemize}

Nada se registra hasta pulsar \textbf{Confirmar}; \textbf{Cancelar} cierra
la ventana sin reservar.

\captura{29c-portal-confirmar-reserva.png}{Confirmación antes de reservar un libro.}

\section{Mis reservas}

\textbf{Mis reservas} (a la derecha en el computador; sobre la lista en el
teléfono) muestra cada reserva con su estado, el plazo para retirarla o tu
lugar en la fila, y permite cancelarla o salir de la fila (también con
confirmación). Después de cancelar o salir de la fila hay que esperar cinco
minutos para pedir otro libro, para que nadie pida y cancele en seguidilla;
mientras tanto, el portal avisa a qué hora puedes volver a pedir.

Cuando tu reserva diga «Lista para retirar», pasa al mesón de préstamo con
tu RUT dentro del plazo indicado (lo recuerda el recuadro \textbf{¿Cómo
retiro un libro?}). Si no la retiras a tiempo, la reserva se cancela sola
esa noche y la copia pasa a la siguiente persona de la fila.

\begin{importante}
Con \textbf{multas pendientes} o con un \textbf{libro vencido} sin
devolver, el portal no te deja reservar ni unirte a la lista de espera: te
indica qué falta (pagar en la biblioteca o devolver el libro).
\end{importante}

% ======================================================================
\chapter{Salas de estudio}
% ======================================================================

Desde el portal reservas las logias y salas de estudio de la biblioteca
\textbf{para el día de hoy}. Antes de reservar, conviene conocer las
reglas:

\begin{itemize}
  \item Las salas funcionan en horario continuo, dentro del horario de la
  biblioteca: de lunes a viernes de 08:45 a 20:30 (la primera reserva
  empieza a las 09:00) y el sábado de 09:00 a 14:00. Los domingos la
  biblioteca no atiende.
  \item Las reservas empiezan a la hora o a la media hora y duran entre 30
  minutos y 2 horas.
  \item Cada reserva es para un grupo de \textbf{3 a 5 personas}, todas
  registradas en la biblioteca (no se reserva para visitantes externos). Los
  seminarios son la excepción: van a nombre de una sola persona, el docente
  o funcionario a cargo.
  \item Cada persona puede sumar hasta \textbf{4 horas de reserva al día},
  no puede estar en dos salas a la vez y tiene \textbf{una sola reserva
  pendiente}: para pedir otra, primero usa o cancela la que tiene.
  \item El grupo tiene \textbf{15 minutos desde la hora de inicio} para
  presentarse en el mesón; si no llega, la sala se libera sola. En fechas de
  mucha demanda, como las de pruebas, la biblioteca puede acortar este tiempo:
  la pantalla de salas siempre muestra el que rige.
\end{itemize}

\section{Ver las salas}

La pantalla muestra cada sala en una \textbf{línea de tiempo} del día: lo
que ya pasó aparece rayado y la línea azul, con la etiqueta \textbf{Ahora},
marca la hora actual. Lo verde está libre. Las reservas de 30 minutos, por
ser cortas, se dibujan como un punto; al pasar el cursor sobre cualquier
bloque se ve su horario. Lo que queda fuera del horario de la biblioteca
también aparece rayado (el sábado, desde las 14:00, con un aviso), y los
domingos la pantalla avisa que la biblioteca no atiende y no permite
reservar.

\captura{30-portal-salas.png}{Reserva de salas de estudio desde el portal, en la línea de tiempo.}

En el teléfono, la línea de tiempo se desliza hacia el lado, con el nombre
de cada sala fijo a la izquierda, y se abre mostrando la hora actual.

Con el botón \textbf{Mapa} ves las salas como están en cada piso de la
biblioteca: toca una para ver sus horarios libres de hoy y reservar desde
ahí. Las salas en gris son las que reserva solo la biblioteca o que
todavía no están en el sistema. El portal recuerda en ese navegador si
preferiste la línea de tiempo o el mapa.

El botón \textbf{¿Cómo reservo una sala?} despliega los pasos, y la página
se actualiza sola cada 30 segundos.

\section{Reservar una sala}

\begin{enumerate}
  \item Toca o haz clic en un espacio verde de la sala que quieras, a la
  hora que quieras.
  \item Elige la duración y escribe el RUT de cada integrante. El tuyo ya
  viene escrito en el primer campo.
  \item Pulsa \textbf{Confirmar}: tu reserva queda en azul en la línea de
  tiempo.
\end{enumerate}

\captura{30b-portal-salas-modal-reserva.png}{Formulario de reserva desde el portal: el primer RUT ya viene con el de tu cuenta.}

Si alguno de los RUT no corresponde a un usuario registrado, o esa persona
ya tiene otra reserva que choca, el formulario marca en rojo exactamente
ese campo.

\section{Cancelar o salir de una reserva}

Tus reservas aparecen en otro color. Al pulsar una que hiciste tú, el
portal pregunta \textbf{¿Cancelar tu reserva?} antes de liberarla. Solo se
cancela una reserva que todavía no se usa: si el grupo ya está en la sala,
se cierra en el mesón al devolver la llave, y si pasó el tiempo para presentarse sin
que llegaran, queda registrada como inasistencia. Las reservas terminadas o
con inasistencia quedan en el historial y no se pueden cancelar.

Puedes cancelar una reserva hecha por ti al día sin consecuencias; desde la
segunda, tienes que esperar cinco minutos para reservar otra sala, y el
portal te lo advierte antes de confirmar la cancelación.

Si otra persona te incluyó en su reserva, no la cancelas para todo el
grupo: al pulsarla, el portal pregunta \textbf{¿Salir de esta reserva?}, y
si confirmas, sales de la lista y el resto del grupo la mantiene. Si en el
grupo quedarían menos de tres personas, la reserva se cancela. Salirte no
te pesa, y aunque estés esperando para reservar, igual puedes ir como
integrante en las reservas de otros.

\section{Llegar y devolver la llave}

Preséntate con tu grupo en el mesón dentro del tiempo que indica la pantalla de
salas (15 minutos, salvo en fechas de mucha demanda) desde la hora
de inicio: el personal confirma la llegada y les entrega la llave. Al
terminar, devuelve la llave en el mesón. Si la devuelves más de 5 minutos
después de la hora de término, el personal puede cursar una multa a nombre
de quien está a cargo de la reserva. Si terminan antes, devolverla no tiene
castigo: el mesón puede liberar el resto del bloque para otro grupo, y
entonces la cuota del día cuenta solo el tiempo que usaron.

% ======================================================================
\chapter{Mi perfil y Mis ingresos}
% ======================================================================

\textbf{Mi perfil} (en el menú de tu cuenta, o \textbf{Perfil} en la barra
inferior del teléfono) muestra tu nombre, tu tipo de usuario, tu RUT, tu
carrera y tu correo, y reúne la \textbf{Constancia de No Multa} (se
bloquea si tienes multas pendientes), el acceso a \textbf{Mis ingresos} y
\textbf{Cerrar sesión}.

\textbf{Mis ingresos} muestra cuántas veces marcaste tu entrada este mes y
tu historial reciente, agrupado por día, con la hora y cómo se registró:
con el código QR, con tu RUT en el portal o en el mesón.

\captura{27c-portal-perfil-ingresos.png}{Mi perfil y Mis ingresos, en el teléfono.}

% ======================================================================
\chapter{Préstamos y multas}
\label{cap:meson}
% ======================================================================

Desde el portal se reserva, pero el préstamo y la devolución se hacen en el
mesón de la biblioteca, con tu RUT (o tu pasaporte o DNI, si eres
estudiante extranjero).

\section{Préstamo de libros}

\begin{itemize}
  \item Cuántos libros puedes llevarte a la vez y por cuántos días depende
  de tu tipo de usuario y del material. El mesón fija la fecha de
  devolución al prestar, y nunca cae en un día en que la biblioteca no
  atiende.
  \item La fecha de devolución de cada préstamo se ve en \textbf{Tu
  actividad}, en el inicio del portal, con un aviso cuando se acerca o ya
  pasó.
  \item Algunos libros se prestan solo para \textbf{uso en sala}: los lees
  dentro de la biblioteca y los devuelves el mismo día, antes del cierre.
  Si no vuelven ese día, quedan atrasados.
  \item Si el libro que buscas está prestado, únete a la lista de espera
  con \textbf{Avisarme al volver}, en el catálogo.
\end{itemize}

\section{Recursos de apoyo}

En el mesón también se prestan recursos de apoyo, como notebooks,
audífonos, cargadores o bicicletas, según lo que tenga la biblioteca. Se
devuelven en el mesón y no generan multa.

\section{Multas}

\begin{itemize}
  \item Si devuelves un libro atrasado, el sistema genera una multa por cada
  día de atraso (en general, \$15 por día, con un tope de \$15.000 por
  préstamo). Los días en que la biblioteca no atiende no se cobran como
  atraso.
  \item También se cursan multas por deterioro o pérdida de material, y el
  personal puede multar la devolución tardía de la llave de una sala.
  \item Las multas se pagan en la biblioteca. Una multa cursada por error se
  anula; ninguna se borra.
\end{itemize}

\begin{importante}
Mientras tengas \textbf{multas pendientes} no puedes llevarte libros ni
recursos de apoyo, reservar libros (en el mesón ni desde el portal) ni
descargar la Constancia de No Multa. Al pagar se desbloquea al instante.
Lo mismo pasa con un \textbf{libro vencido}: hasta que lo devuelvas, no
puedes llevarte ni reservar más material.
\end{importante}

% ======================================================================
\chapter{Preguntas Frecuentes}
% ======================================================================

\begin{description}[leftmargin=0pt,style=nextline]
  \item[Escribo mi RUT y no puedo entrar. ¿Qué pasa?]
  Al portal se entra con tu \textbf{correo institucional} (por ejemplo,
  \texttt{nombre@umag.cl}) y tu contraseña, no con el RUT.

  \item[Olvidé mi contraseña. ¿Qué hago?]
  Acércate al mesón de la biblioteca con tu credencial: el personal te
  asigna una nueva. El portal no tiene recuperación en línea.

  \item[No puedo entrar al portal ni marcar mi entrada.]
  Revisa que estés usando tu correo. Si el problema sigue, puede que tu
  cuenta haya vencido (vence el 31 de marzo): acércate al mesón para
  renovarla.

  \item[¿Por qué no puedo reservar un libro?]
  Porque tienes una multa pendiente o un libro vencido sin devolver, o
  porque cancelaste otra reserva hace menos de cinco minutos. El portal te
  dice cuál es el motivo.

  \item[¿Puedo reservar una sala para la próxima semana?]
  No. Desde el portal solo se reserva para el día de hoy (y nunca un
  domingo, que la biblioteca no atiende).

  \item[¿Qué pasa si reservamos una sala y no llegamos?]
  El grupo tiene 15 minutos desde la hora de inicio para presentarse en el
  mesón, o el tiempo que indique la pantalla de salas si la biblioteca lo
  cambió, como en fechas de pruebas. Si no llega, la sala se libera sola y queda
  registrada la inasistencia.

  \item[Me incluyeron en una reserva de sala sin preguntarme.]
  Pulsa la reserva y elige \textbf{Salir de esta reserva}. El resto del
  grupo la mantiene, y a ti no te pesa.

  \item[¿Puedo marcar la entrada de un compañero?]
  No. El RUT que se escribe en Marcar Entrada tiene que ser el de la cuenta
  con que iniciaste sesión.

  \item[Aparece el aviso ``No se pudo conectar con el servidor''.]
  Mientras no vuelva la conexión, nada de lo que se ve en pantalla está al
  día. Revisa tu conexión a internet y vuelve a intentarlo.
\end{description}

% ======================================================================
\chapter*{Glosario}
\addcontentsline{toc}{chapter}{Glosario}
% ======================================================================

\begin{description}[leftmargin=0pt,style=nextline]
  \item[Ejemplar] Una copia física concreta de un libro, con su propio
  código de barras. Un mismo título puede tener varios ejemplares.

  \item[Lista de espera] Personas esperando que se libere una copia de un
  libro, en orden de llegada: la primera en pedirlo es la primera en
  recibirlo.

  \item[Logia] Nombre que usa la biblioteca para sus salas de estudio
  grupales, que se reservan en horario continuo por hasta dos horas cada
  vez.

  \item[Multa] Deuda con la biblioteca. Las de atraso se generan solas al
  devolver un libro fuera de plazo; las de deterioro, pérdida u otro motivo
  las cursa el personal. Se pagan o se anulan, pero nunca se borran.

  \item[Constancia de No Multa] Documento en PDF que certifica que no
  tienes multas pendientes con la biblioteca. Se descarga desde el portal.

  \item[Portal de la biblioteca] El sitio de autoservicio para que
  estudiantes, docentes y funcionarios resuelvan trámites simples por su
  cuenta, sin pasar por el mesón.
\end{description}

\end{document}
